\subsection{Класс DHeap}

Класс $DHeap$ реализует $d$-арную кучу, где $d$ — количество дочерних узлов для каждого узла.

\begin{itemize}[label=•, labelsep=1em]
    \item \texttt{Конструктор} принимает целое значение $d$ и проверяет, что оно больше или равно 2. Если $d < 2$, генерируется исключение.
\end{itemize}

\begin{itemize}[label=•, labelsep=1em]
    \item \texttt{Insert} вставляет новый элемент в кучу. Элемент добавляется в конец вектора, после чего вызывается функция \texttt{SiftUp} для поддержания свойств кучи.
\end{itemize}

\begin{itemize}[label=•, labelsep=1em]
    \item \texttt{ExtractMin} извлекает минимальный элемент из кучи. Если куча пуста, генерируется исключение. Минимальный элемент сохраняется, последний элемент вектора перемещается на место корня, после чего вызывается \texttt{SiftDown}.
\end{itemize}
\begin{lstlisting}[language=Pascal]
    function ExtractMin(var v: array of Integer): Integer;
    begin
        if IsEmpty(v) then
            raise Exception.Create('Heap is empty');
    
        var min_element := v[0];
        v[0] := v[High(v)];
        SetLength(v, Length(v) - 1);
        if not IsEmpty(v) then
            SiftDown(v, 0);
    
        Result := min_element;
    end;
\end{lstlisting}

\begin{itemize}[label=•, labelsep=1em]
    \item \texttt{IsEmpty} проверяет, пуста ли куча.
\end{itemize}

\begin{itemize}[label=•, labelsep=1em]
    \item \texttt{SiftDown} восстанавливает свойства кучи, начиная с узла $i$. Перемещает элемент вниз по дереву, пока он не станет больше своих дочерних элементов.
\end{itemize}
\begin{lstlisting}[language=Pascal]
    procedure SiftDown(var v: array of Integer; index: Integer);
    var
        smallest, left, right: Integer;
    begin
        smallest := index;
        left := 2 * index + 1;
        right := 2 * index + 2;
    
        if (left < Length(v)) and (v[left] < v[smallest]) then
            smallest := left;
    
        if (right < Length(v)) and (v[right] < v[smallest]) then
            smallest := right;
    
        if smallest <> index then
        begin
            var temp := v[index];
            v[index] := v[smallest];
            v[smallest] := temp;
    
            SiftDown(v, smallest);
        end;
    end;
    \end{lstlisting}

    \begin{itemize}[label=•, labelsep=1em]
        \item \texttt{SiftUp} восстанавливает свойства кучи, начиная с узла $i$. Перемещает элемент вверх по дереву, пока он не станет меньше своих родительских узлов.
    \end{itemize}
    
    \begin{lstlisting}[language=Pascal]
    procedure DHeap.SiftUp(i: Integer);
    var
        key0, value0: Integer;
        p: Integer;
    begin
        key0 := v_[i].first;
        value0 := v_[i].second;
        p := Father(i);
    
        while (i <> 0) and (v_[p].second > value0) do
        begin
            v_[i] := v_[p];
            i := p;
            p := Father(i);
        end;
        v_[i] := (key0, value0);
    end;
    \end{lstlisting}
    
    \begin{itemize}[label=•, labelsep=1em]
        \item \texttt{MinChild} находит индекс минимального дочернего элемента для узла $i$. Если у узла нет дочерних элементов, возвращается индекс самого узла.
    \end{itemize}
    
    \begin{itemize}[label=•, labelsep=1em]
        \item \texttt{LastChild} возвращает индекс последнего дочернего элемента для узла $i$. Если у узла нет дочерних элементов, возвращается -1.
    \end{itemize}
    
    \begin{itemize}[label=•, labelsep=1em]
        \item \texttt{FirstChild} возвращает индекс первого дочернего элемента для узла $i$. Если узел не имеет дочерних элементов, возвращается -1.
    \end{itemize}
    
    \begin{itemize}[label=•, labelsep=1em]
        \item \texttt{Father} возвращает индекс родительского узла для узла $i$.
    \end{itemize}